\section{A minimum exists}\label{sec:existence}
We take a convex body to mean a compact convex subset of $\mathbb R^3$
with nonempty interior. Let $\mathcal U$ be the class of bodies containing,
for each rectifiable curve of length at most one, a copy obtained by a
translation and a proper rotation. The infimum
\[
 V_3 = \inf_{K\in\mathcal U}|K|
\]
is a minimum.

First, $\mathcal U$ is nonempty: every unit-speed curve on $[0,1]$ is
within distance $1/2$ of its point at parameter $1/2$, so the ball of radius
$1/2$ is universal. Shorter curves can be extended to length one.
Next consider the polygonal curve that traverses three successive edges
of a regular tetrahedron $T$ of edge length $a=1/3$. It has length one
and visits all four vertices, so its convex hull is $T$. This tetrahedron
has volume $a^3/(6\sqrt2)$ and surface area $\sqrt3a^2$; consequently
its insphere has radius
\[
 r_0=\frac{3|T|}{\operatorname{area}(\partial T)}
 =\frac{a\sqrt6}{12}=\frac{\sqrt6}{36}>0.
\]
Every universal body contains a congruent copy of $T$.

Take a minimizing sequence $K_j$ with $|K_j|\le M$ for some finite $M$.
Apply a proper rigid motion separately to each $K_j$ to make its forced
copy of $T$ equal to a single fixed $T$, translated so that the incenter
is the origin. This preserves universality and volume. Thus all $K_j$
contain $B(0,r_0)$. For any nonzero point $x\in K_j$, the convex hull of $x$ and
the disk $B(0,r_0)\cap x^\perp$ is a cone contained in $K_j$. Its volume is
$\pi r_0^2|x|/3$, and therefore
\[
 |x|\le R:=\frac{3M}{\pi r_0^2}.
\]
The bodies are now uniformly bounded and all contain the same ball.

By the Blaschke selection theorem there is a Hausdorff-convergent
subsequence $K_j\to K$; its limit is compact and convex and contains
$B(0,r_0)$. For completeness, the compactness statement here can also be
obtained by applying Arzela--Ascoli to the support functions
$u\mapsto\max_{x\in K_j}u\cdot x$, which are uniformly bounded by $R$ and
are $R$-Lipschitz on the unit sphere.

We check universality directly. Fix any curve $\gamma$, translate its
parametrization initially so that $\gamma(0)=0$, and choose rotations
$Q_j\in SO(3)$ and vectors $b_j$ with
$Q_j\gamma+b_j\subseteq K_j$. Then $b_j\in K_j\subseteq B(0,R)$.
Compactness of $SO(3)\times B(0,R)$ gives, after a further subsequence
which is permitted to depend on this fixed curve,
$(Q_j,b_j)\to(Q,b)$. The compactness of the curve image and Hausdorff
convergence imply $Q\gamma+b\subseteq K$. This proves universality for
each curve; a common subsequence of placement maps for all curves is
neither claimed nor needed.

Finally, volume converges. If $\delta_j=d_H(K_j,K)$, then
\[
 K_j\subseteq K+\delta_j B(0,1)
 \subseteq(1+\delta_j/r_0)K,
\]
and the same inclusion holds with $K$ and $K_j$ interchanged. Taking
volumes and letting $\delta_j\to0$ gives $|K_j|\to|K|$. Hence
$|K|=V_3$, as asserted.

